\documentclass[11pt,a4paper]{article}
\usepackage[utf8]{inputenc}
\usepackage[T1]{fontenc}
\usepackage{amsmath,amssymb,amsthm}
\usepackage{hyperref}

\newtheorem{theorem}{Theorem}
\newtheorem{remark}{Remark}

\title{\textbf{Pilar 86 Quantum Kernel Architecture and $\phi$-Aware Distributed Operating System: Leech Lattice ($\Lambda_{24}$) Compilation, Non-Commutative Logic Gates, and Reversible Transition}}
\author{Mariano Panzano Caball\'e \\ \small Oasis Sovereign Architecture Research, Barcelona}
\date{November 17, 2026}

\begin{document}
\maketitle

\begin{abstract}
Building upon our machine-checked Lean 4 verification series (DOIs: \href{https://doi.org/10.5281/zenodo.24311596}{10.5281/zenodo.24311596} and \href{https://doi.org/10.5281/zenodo.24323807}{10.5281/zenodo.24323807}), we present the formal specification of the \textbf{Pilar 86 Quantum Kernel} and the five-layer \textbf{$\phi$-Aware Distributed Operating System}. To bridge contemporary CMOS/FPGA classical hosts with fault-tolerant Quantum Processing Units (QPUs), the architecture integrates: (i) critical damping under the decadic attractor $\kappa = \ln 10 \approx 2.302585$; (ii) logical gate compilation mapped onto the 24-dimensional Leech lattice ($\Lambda_{24}$) and Monster symmetry representations ($196883 = 47 \times 59 \times 71$); (iii) a structural phase circuit-breaker anchored on Prime 41 to preserve state purity ($\mathrm{Tr}(\rho^2) = 1$); (iv) a minimal 5-generator WASM kernel (\texttt{LieAlgebraCore}, \texttt{BlochStateEngine}) paired with non-commutative \texttt{QuantumLogicGate} security ($A \cdot B \neq B \cdot A$) and ZX-Calculus optimization; and (v) Fibonacci state-space pruning ($\phi^N < 2^N$). Structural bounds are machine-checked in Lean 4 (\texttt{Pilar86\_LeechQPU\_QuantumOS\_Fase11.lean}) with zero unproven axioms.
\end{abstract}

\section{Pilar 86 QPU Kernel Primitives}
\begin{enumerate}
    \item \textbf{Leech Lattice ($\Lambda_{24}$) Compilation:} Embedding multi-qubit entanglement graphs into the densest 24-dimensional sphere packing minimizes Euclidean code distance overhead and absorbs channel deformations geometrically.
    \item \textbf{Prime-41 Phase Circuit-Breaker:} Monitors purity drift ($\mathrm{Tr}(\rho^2) < 1$) and flushes noisy register accumulations deterministically before decoherence cascades.
    \item \textbf{Lie-Algebraic Control \& Non-Commutative Security:} Continuous resource scheduling via $\mathfrak{su}(2)/\mathfrak{su}(n)$ infinitesimal generators and Euler rotations $R_z(\theta) = e^{-i\theta\sigma_z/2}$ replaces abrupt binary interrupts, while non-commutative projector ordering enforces temporal race-condition immunity.
\end{enumerate}

\section{Machine-Checked Structural Theorems (Lean 4)}
\begin{theorem}[Fibonacci State-Space Pruning Bound]
For every register depth $n \in \mathbb{N}$, the Fibonacci-constrained transition cardinality satisfies $F_n \le 2^n$, strictly bounding exponential state explosion in $\phi$-aware schedulers.
\end{theorem}

\section{Experimental Boundaries and Physical Validation Roadmap}
In strict adherence to thermodynamic rigor, classical CMOS/WASM simulations and P2P Langevin virtual quantum overlays do not constitute physical QPU validation nor refute Landauer's classical erasure bound ($k_B T \ln 2$) in thermal equilibrium. Empirical verification of the adiabatic Fibonacci dissipation regime ($k_B T \ln \phi$) requires direct calorimetric and coherence measurements on physical superconducting qubits and Fibonacci anyon platforms.

\begin{thebibliography}{9}
\bibitem{panzano2026a} M. Panzano Caball\'e, \emph{Oasis Sovereign Monolith: Machine-Checked Lean 4 Spectral Bounds in Discretized SU(2) Gauge Models and Decadic Entropy Confinement}, Zenodo, 2026. \url{https://doi.org/10.5281/zenodo.24311596}
\bibitem{panzano2026c} M. Panzano Caball\'e, \emph{Araki Relative Entropy Monotonicity and Variational Matrix Compression in Machine-Checked Cold-Silicon Gauge Verification}, Zenodo, 2026. \url{https://doi.org/10.5281/zenodo.24323807}
\bibitem{qec_noise2026} \emph{Resolving and Suppressing Temporally Correlated Noise in Quantum Error Correction}, arXiv:2611.21578v1 [quant-ph], Nov. 2026.
\bibitem{verify_qec2026} \emph{Specifying \& Verifying Quantum Error-Correcting Codes}, arXiv:2611.21236v1 [quant-ph], Nov. 2026.
\end{thebibliography}

\end{document}
